\documentclass{beamer}

\usepackage{amsmath}

\begin{document}

\begin{frame}
	\frametitle{Pairwise correlations among Vancouver vs suburban types}
	\begin{tabular}{rrrrr}
  \hline
 & VanSingle & VanCondo & FVSingle & FVCondo \\
  \hline
VanSingle & 1.00 & 0.71 & 0.82 & 0.54 \\
  VanCondo & 0.71 & 1.00 & 0.71 & 0.70 \\
  FVSingle & 0.82 & 0.71 & 1.00 & 0.71 \\
  FVCondo & 0.54 & 0.70 & 0.71 & 1.00 \\
   \hline
\end{tabular}
\end{frame}

\begin{frame}
	\frametitle{Useful reference? Nihilism vs hedonism?}
	\framesubtitle{Berliant/McMillen}
	\begin{itemize}
		\item Might make same linear gradient choice?
	\end{itemize}
\end{frame}

\begin{frame}
	\frametitle{So what would be helpful in intro/conclusion}
	\framesubtitle{Some ad hoc thougths}
	\begin{itemize}
		\item What decision would be affected by different index choice?
		\item Portfolio efficiency for investor?
		\item Bailout decision for defaulters?
		\item Rationalize cap rate differences?
	\end{itemize}
\end{frame}

\begin{frame}
	\frametitle{Maybe overly bold claim}
	\begin{itemize}
		\item ``insufficient attention has been paid to the implications of non-random sampling within cities''
		\item Many people were thinking about defaults polluting indexes (or lack thereof biasing upward) during GFC
		\item If no pubs maybe because no one found anything?
		\item Vague recollection I tried and found nothing
		\item But location of course within metro.
		\item See my review of Handbook of Urban Econ thing \dots
	\end{itemize}
\end{frame}

\begin{frame}
	\frametitle{Zip code price movements boom bust (source, Zillow)}
	\framesubtitle{Red: Flyover, Blue: Coastal, green: Sand State}
	\pgfimage[width=.75\textwidth]{zips}
\end{frame}

\begin{frame}
	\frametitle{Price of services}
	``If the objective is to measure the price of housing services or
	housing space, then the definitions of hand $H$ in the preceding
	equations should reflect a measure of housing services, rather than
	assuming that each housing unit provides the same quantity of housing
	space.''
	\begin{itemize}
		\item Is the point that Bill Gates's house should get extra weight?
	\end{itemize}
\end{frame}

\begin{frame}
	\frametitle{Muth discussion}
	\begin{itemize}
		\item What is going on with $\tau k$ and $\frac{\partial \tau}{\partial k}$ Linear? Nonlinear?
		\item Per above, linear in distance not realistic due to congestion. Going 1 mile on Cross-Bronx Expressway slower than one mile on 95 in the suburbs.
	\end{itemize}
\end{frame}

\begin{frame}
	\frametitle{Common appreciation}
	\begin{itemize}
		\item Standard monocentric gradient linear
			\begin{itemize}
				\item Value at $r$ = structure + lot size $\theta \int_{s=r}^{b} \frac{1}{\text{lot size at s}}ds$
				\item Time log derivative unlikely constant by $r$
			\end{itemize}
		\item But no less realistically (per Williams) $\log v = \log \text{structure} + \theta \left[\ln b(t) - \ln r\right]$
			\begin{itemize}
				\item Then location not clear a source of growth heterogeneity
				\item Might be fine with, say unit income elastcity
			\end{itemize}
		\item Zoned single family vs multifamily clearer?
	\end{itemize}
\end{frame}

\begin{frame}
	\frametitle{Proposition 2 doesn't really follow from Prop 1}
	\begin{itemize}
		\item Are you missing ``if not quasi-linearity, and if linear gradient in distance''?
	\end{itemize}
\end{frame}

\begin{frame}
	\frametitle{Is AHS price data bananas?}
	\begin{itemize}
		\item 2015-2019 Table 1 shows 18\% vs 13\% per year appreciation
		\item FHFA per FRED says about 7\%
		\item Washington, D.C. Table 1 says about 2\%
		\item FHFA per FRED says about 8\%.
		\item Even I can afford Attom \dots
		\item Aggregate Zillow zip indexes by metro central/suburbs?
		\item Given giant idiosyncratic errors, what can we rule out?
		\item Strong suggestion: use better data and see if differences as large
	\end{itemize}
\end{frame}

\end{document}

